% META: {"id": "TEXP-GRA-EX-005", "chapitre": "TEXP-GRAPHES", "type_objet": "exercice", "capacites": ["TEXP-GRA-C5"], "capacites_codes": ["C5"], "methodes": [], "parcours": 3, "competences": ["raisonner", "calculer"], "duree_min": 18, "mode_creation": "ex_nihilo", "sources_inspiration": [], "fichier_tex": "chapitres/TEXP-GRAPHES/exercices/TEXP-GRA-EX-005.tex", "corrige_tex": "chapitres/TEXP-GRAPHES/corriges/TEXP-GRA-CO-005.tex", "status": "generated"}
% BEGIN-VERIFY
% from sympy import Matrix
% M = Matrix([
%     [0,1,1,0],
%     [1,0,1,0],
%     [1,1,0,1],
%     [0,0,1,0],
% ])
% M2 = M**2
% M3 = M**3
% # D (indice 3) n'a qu'un seul voisin : C (indice 2)
% voisins_D = [k for k in range(4) if M[k,3] == 1]
% assert voisins_D == [2]
% somme_recurrence = sum(M2[0, k] * M[k, 3] for k in range(4))
% assert somme_recurrence == M3[0, 3]
% assert somme_recurrence == M2[0, 2] * M[2, 3]
% assert M2[0,2] == 1
% END-VERIFY
\begin{exercice}{TEXP-GRA-EX-005}{3}{18}
On reprend le graphe $H$ (sommets $A,B,C,D$) et ses matrices $M$, $M^2$, $M^3$ de l'exercice précédent.

On rappelle la relation de récurrence démontrée en cours : pour tous sommets $i,j$,
$$(M^{n+1})_{ij}=\sum_{k}(M^n)_{ik}\times m_{kj}$$
où la somme porte en réalité uniquement sur les sommets $k$ adjacents à $j$.
\begin{enumerate}
  \item Quel est l'unique voisin du sommet $D$ dans $H$ ?
  \item En déduire que $(M^3)_{A,D}=(M^2)_{A,C}$, sans recalculer entièrement le produit $M^2\times M$.
  \item Sachant que $(M^2)_{A,C}=1$, retrouver la valeur de $(M^3)_{A,D}$ trouvée à l'exercice précédent.
  \item Expliquer, avec les mots du cours (décomposition d'une chaîne), pourquoi ce raisonnement est valable.
\end{enumerate}
\end{exercice}
